\documentclass[11pt,letterpaper]{article}
\PassOptionsToPackage{hyphens}{url}
\usepackage[margin=.83in,top=.8in,bottom=.8in,headheight=15pt]{geometry}
\usepackage[T1]{fontenc}
\usepackage{lmodern,amsmath,amssymb,mathtools,bm,booktabs,array,graphicx,tikz}
\usetikzlibrary{arrows.meta,positioning,calc,decorations.pathreplacing,patterns}
\definecolor{ink}{HTML}{17344D}\definecolor{accent}{HTML}{087E8B}\definecolor{pale}{HTML}{EDF5F5}\definecolor{gold}{HTML}{AD6B19}
\providecommand{\dd}{\,\mathrm{d}}\providecommand{\R}{\mathbb{R}}
\providecommand{\E}{\mathbb{E}}\providecommand{\Var}{\operatorname{Var}}
\providecommand{\Cov}{\operatorname{Cov}}\providecommand{\diag}{\operatorname{diag}}
\providecommand{\norm}[1]{\left\lVert#1\right\rVert}
\providecommand{\abs}[1]{\left|#1\right|}
\usepackage{microtype,hyperref,fancyhdr,enumitem,amsthm,sectsty,needspace}
\allsectionsfont{\raggedright}
\hypersetup{colorlinks=true,linkcolor=ink,urlcolor=accent,pdfauthor={Zan Ahmad},pdftitle={Crossbridge populations: transition kinetics and mechanical force / Review}}
\setlength{\parindent}{0pt}\setlength{\parskip}{6pt}\linespread{1.07}
\setlist{itemsep=3pt,topsep=4pt,leftmargin=*}
\emergencystretch=2em\clubpenalty=10000\widowpenalty=10000
\pagestyle{fancy}\fancyhf{}\fancyhead[L]{\small Mathematical Modeling}\fancyhead[R]{\small Lecture 5}
\fancyfoot[L]{\footnotesize Zan Ahmad}\fancyfoot[R]{\footnotesize\thepage}\renewcommand{\headrulewidth}{.3pt}
\newtheorem{proposition}{Proposition}\newtheorem{theorem}{Theorem}\newtheorem{lemma}{Lemma}
\theoremstyle{definition}\newtheorem{definition}{Definition}\newtheorem{example}{Example}
\numberwithin{equation}{section}
\begin{document}
\section*{Crossbridge populations: transition kinetics and mechanical force --- review sheet}
Use this after attempting the worked examples and practice. Explain each claim in words before calculating.

\section*{Key ideas}
\begin{itemize}
\item Attachment and detachment fluxes depend on the available source fractions.
\item Equal opposing fluxes give a steady population fraction without stopping individual transitions.
\item For positive rate sum, equilibrium is f/(f+\allowbreak{}g) and relaxation time is 1/(f+\allowbreak{}g).
\item A fixed-force readout adds assumptions that occupancy alone does not establish.
\end{itemize}
\section*{Common errors and repairs}
\textbf{Treating a rate constant as a fraction or a total flux.} Rates have units 1/s; multiply by the available fraction to obtain a transition flux.

\textbf{Doubling both rates must double the final fraction.} The ratio is unchanged, so the equilibrium is unchanged; the approach is faster.

\textbf{A force trace validates a force-velocity theory.} The model does not contain velocity or strain, so that claim exceeds its equations.

\section*{Retrieve without your notes}
\begin{itemize}
\item At f =\allowbreak{} 3/s, g =\allowbreak{} 1/s and A =\allowbreak{} 0.75, what are the two fluxes?
\item What is carried across an instantaneous change in rate constants?
\item Name one independent observation needed to test the occupancy explanation.
\end{itemize}
\section*{Readiness check}
\begin{itemize}
\item Separate a state fraction from an attachment or detachment rate.
\item Compute the competing transition fluxes and predict whether attachment increases.
\item Find the equilibrium fraction by balancing gains and losses.
\item Explain why scaling both rates changes the timescale but not the equilibrium.
\item Identify what additional information is needed to predict shortening or metabolic energy.
\end{itemize}
\end{document}
